\section{Data Availability (DA)}

\label{sec:data_availability}

Rollup proofs verify computational correctness, but actual transaction data (or state diffs) is required to independently reconstruct and verify the rollup state. This concept is termed Data Availability (DA) \citep{saif2024dataavailability}. For a rollup to be secure and trustless, the sequencer must publish this data so that anyone can run an independent node. If a rollup operator withholds data, users cannot verify the state or recover funds. It must be noted that DA is not simply 'metadata'—it is the core underlying transaction data necessary to derive the state transition.

Historically, ZK-Rollups posted this data to Ethereum L1 using standard transaction \texttt{calldata}. However, \texttt{calldata} is permanently stored on all Ethereum full nodes, making it prohibitively expensive and typically the dominant component of L2 transaction costs.

To mitigate this, Ethereum introduced EIP-4844 ("Proto-Danksharding"), which created a new transaction type that carries "blobs" of data \citep{eip4844}. Blobs are temporarily available on the consensus layer (typically for a few weeks) and are priced in a separate, significantly cheaper fee market. While the EVM execution layer cannot read the blob data directly, it can verify cryptographic commitments (such as KZG commitments) to the blob data \citep{eip4844}. This fundamental shift significantly alters the cost structure of DA for rollups \citep{saif2024dataavailability, lavaur2023modular}. Off-chain DA solutions offer even cheaper alternatives but introduce different trust assumptions, pushing the architecture toward a Validium model rather than a pure rollup.
